
\documentclass{article}
\usepackage{colortbl}
\usepackage{makecell}
\usepackage{multirow}
\usepackage{supertabular}

\begin{document}

\newcounter{utterance}

\centering \large Interaction Transcript for game `referencegame', experiment `line\_grids\_columns', episode 13 with reflexion{-}llama{-}3.1{-}8b{-}instruct.
\vspace{24pt}

{ \footnotesize  \setcounter{utterance}{1}
\setlength{\tabcolsep}{0pt}
\begin{supertabular}{c@{$\;$}|p{.15\linewidth}@{}p{.15\linewidth}p{.15\linewidth}p{.15\linewidth}p{.15\linewidth}p{.15\linewidth}}
    \# & $\;$A & \multicolumn{4}{c}{Game Master} & $\;\:$B\\
    \hline

    \theutterance \stepcounter{utterance}  
    & & \multicolumn{4}{p{0.6\linewidth}}{
        \cellcolor[rgb]{0.9,0.9,0.9}{
            \makecell[{{p{\linewidth}}}]{
                \texttt{\tiny{[P1$\langle$GM]}}
                \texttt{You have attempted this same task before and did not succeed. These are the plans you wrote after those attempts.} \\
\texttt{Trial \#1: New plan: I failed to distinguish the target grid from the distractor grids because my description focused on the overall pattern of X's, which is present in all three grids. What I should have done is to focus on a unique feature of the target grid that is not present in the distractor grids. In this case, I should have noticed that the target grid has a unique pattern of X's in the first and last columns, where each row has an X in the first column and an X in the last column. Next time, I will focus on this unique feature to generate a referring expression that accurately describes the target grid and distinguishes it from the distractor grids.} \\
\\ 
\texttt{Trial \#2: New plan: I failed to distinguish the target grid from the distractor grids because my description focused on the overall pattern of X's, which is present in all three grids. What I should have done is to focus on a unique feature of the target grid that is not present in the distractor grids. In this case, I should have noticed that the target grid has a unique pattern of X's in the first and last columns, where each row has an X in the first column and an X in the last column. However, I also failed to notice that the target grid has a unique pattern of X's in the middle column, where each row has two X's. Next time, I will focus on this unique feature to generate a referring expression that accurately describes the target grid and distinguishes it from the distractor grids. I will also make sure to carefully examine the distractor grids to ensure that my description is not applicable to them.} \\
\\ 
\texttt{Trial \#3: New plan: I failed to distinguish the target grid from the distractor grids because my description focused on the overall pattern of X's, which is present in all three grids. What I should have done is to focus on a unique feature of the target grid that is not present in the distractor grids. In this case, I should have noticed that the target grid has a unique pattern of X's in the first and last columns, where each row has an X in the first column and an X in the last column. However, I also failed to notice that the target grid has a unique pattern of X's in the middle column, where each row has two X's. Next time, I will focus on this unique feature to generate a referring expression that accurately describes the target grid and distinguishes it from the distractor grids. I will also make sure to carefully examine the distractor grids to ensure that my description is not applicable to them. Additionally, I will try to identify a combination of features that are unique to the target grid, rather than relying on a single feature. For example, I could describe the target grid as the 5x5 grid with X's in the first and last columns, and two X's in the middle column, which is not present in the distractor grids.} \\
\\ 
\\ 
\texttt{You are given three grids, where each of them is 5 by 5 in size.} \\
\texttt{Grids have empty cells marked with "▢" and filled cells marked with "X".} \\
\texttt{Your task is to generate a referring expression that best describes the target grid while distinguishing it from the two other distractor grids.} \\
\texttt{The first grid is the target grid, and the following two grids are the distractors.} \\
\\ 
\texttt{Target grid:} \\
\\ 
\texttt{□ X X □ X} \\
\texttt{□ X X □ X} \\
\texttt{□ X X □ X} \\
\texttt{□ X X □ X} \\
\texttt{□ X X □ X} \\
\\ 
\texttt{Distractor grid 1:} \\
\\ 
\texttt{□ □ X X X} \\
\texttt{□ □ X X X} \\
\texttt{□ □ X X X} \\
\texttt{□ □ X X X} \\
\texttt{□ □ X X X} \\
\\ 
\texttt{Distractor grid 2:} \\
\\ 
\texttt{□ X X X □} \\
\texttt{□ X X X □} \\
\texttt{□ X X X □} \\
\texttt{□ X X X □} \\
\texttt{□ X X X □} \\
\\ 
\texttt{Instruction: Describe the target grid.} \\
\texttt{Generate the referring expression starting with the tag "Expression: " for the given target grid. Omit any other text.} \\
            }
        }
    }
    & \\ \\

    \theutterance \stepcounter{utterance}  
    & \multicolumn{4}{p{0.6\linewidth}}{
        \cellcolor[rgb]{0.9,0.9,0.9}{
            \makecell[{{p{\linewidth}}}]{
                \texttt{\tiny{[P1$\rangle$GM]}}
                \texttt{Expression: The 5x5 grid with X's in the first and last columns, and two X's in each of the middle three columns.} \\
            }
        }
    }
    & & \\ \\

    \theutterance \stepcounter{utterance}  
    & & & \multicolumn{2}{p{0.3\linewidth}}{
        \cellcolor[rgb]{0.9,0.9,0.9}{
            \makecell[{{p{\linewidth}}}]{
                \texttt{\tiny{[GM$|$GM]}}
                \texttt{Expression: The 5x5 grid with X's in the first and last columns, and two X's in each of the middle three columns.} \\
            }
        }
    }
    & & \\ \\

    \theutterance \stepcounter{utterance}  
    & & \multicolumn{4}{p{0.6\linewidth}}{
        \cellcolor[rgb]{0.9,0.9,0.9}{
            \makecell[{{p{\linewidth}}}]{
                \texttt{\tiny{[GM$\rangle$P2]}}
                \texttt{You have attempted this same task before and did not succeed. These are the plans you wrote after those attempts.} \\
\texttt{Trial \#1: New plan: I failed to match the given expression because I didn't carefully analyze the description. The expression states that the grid has four rows of five X's, with the first and last X in each row being in the same column. However, I selected the third grid, which has five rows of five X's, not four. Next time, I will carefully read the expression and look for a grid that matches the description. I will pay attention to the number of rows and the position of the X's in each row.} \\
\\ 
\texttt{Trial \#2: New plan: I failed to match the given expression because I didn't carefully analyze the description. The expression states that the grid has X's in the first and last columns of each row. However, I selected the third grid, which has X's in the first and last columns of each row, but not in the first and last columns of each row as the expression states. The correct interpretation is that the first and last X in each row are in the same column, not necessarily the first and last columns. Next time, I will carefully read the expression and look for a grid that matches the description. I will pay attention to the position of the X's in each row and the column in which the first and last X appear.} \\
\\ 
\texttt{Trial \#3: New plan: I failed to match the given expression because I didn't carefully analyze the description. The expression states that the grid has X's in the first and last columns of each row, but I selected the third grid, which has X's in the first and last columns of each row, but not in the first and last columns of each row as the expression states. The correct interpretation is that the first and last X in each row are in the same column, not necessarily the first and last columns. However, I also failed to notice that the expression states "two X's in the middle column", which is not present in the third grid. Next time, I will carefully read the expression and look for a grid that matches the description. I will pay attention to the position of the X's in each row, the column in which the first and last X appear, and the presence of X's in the middle column.} \\
\\ 
\\ 
\texttt{You are given three grids, where each of them is 5 by 5 in size.} \\
\texttt{Grids have empty cells marked with "▢" and filled cells marked with "X".} \\
\texttt{You are also given a referring expression that describes one of the given grids.} \\
\texttt{Your task is to select the grid that matches the given referring expression.} \\
\texttt{Generate only the number (in text) of the grid that matches the given expression by selecting first, second, or third.} \\
\\ 
\texttt{First:} \\
\\ 
\texttt{□ □ X X X} \\
\texttt{□ □ X X X} \\
\texttt{□ □ X X X} \\
\texttt{□ □ X X X} \\
\texttt{□ □ X X X} \\
\\ 
\texttt{Second:} \\
\\ 
\texttt{□ X X □ X} \\
\texttt{□ X X □ X} \\
\texttt{□ X X □ X} \\
\texttt{□ X X □ X} \\
\texttt{□ X X □ X} \\
\\ 
\texttt{Third:} \\
\\ 
\texttt{□ X X X □} \\
\texttt{□ X X X □} \\
\texttt{□ X X X □} \\
\texttt{□ X X X □} \\
\texttt{□ X X X □} \\
\\ 
\texttt{Expression: The 5x5 grid with X's in the first and last columns, and two X's in each of the middle three columns.} \\
\texttt{Question: Which grid does the expression refer to?} \\
\texttt{Start with the tag "Answer: ", followed by your selection. Omit any other text.} \\
            }
        }
    }
    & \\ \\

    \theutterance \stepcounter{utterance}  
    & & & \multicolumn{4}{p{0.6\linewidth}}{
        \cellcolor[rgb]{0.9,0.9,0.9}{
            \makecell[{{p{\linewidth}}}]{
                \texttt{\tiny{[GM$\langle$P2]}}
                \texttt{Answer: 3} \\
            }
        }
    }
     \\ \\

    \theutterance \stepcounter{utterance}  
    & & & \multicolumn{2}{p{0.3\linewidth}}{
        \cellcolor[rgb]{0.9,0.9,0.9}{
            \makecell[{{p{\linewidth}}}]{
                \texttt{\tiny{[GM$|$GM]}}
                \texttt{3} \\
            }
        }
    }
    & & \\ \\

    \theutterance \stepcounter{utterance}  
    & \multicolumn{4}{p{0.6\linewidth}}{
        \cellcolor[rgb]{0.95,0.95,0.95}{
            \makecell[{{p{\linewidth}}}]{
                \texttt{\tiny{[P2$\rightarrow$P1]}}
                \texttt{Post{-}episode Reflexion (not sent to GM):} \\
\texttt{New plan: I failed to distinguish the target grid from the distractor grids because my description focused on the overall pattern of X's, which is present in all three grids. What I should have done is to focus on a unique feature of the target grid that is not present in the distractor grids. In this case, I should have noticed that the target grid has a unique pattern of X's in the first and last columns, where each row has an X in the first column and an X in the last column. However, I also failed to notice that the target grid has a unique pattern of X's in the middle column, where each row has two X's. Next time, I will focus on this unique feature to generate a referring expression that accurately describes the target grid and distinguishes it from the distractor grids. I will also make sure to carefully examine the distractor grids to ensure that my description is not applicable to them. Additionally, I will try to identify a combination of features that are unique to the target grid, rather than relying on a single feature. For example, I could describe the target grid as the 5x5 grid with X's in the first and last columns, and two X's in the middle column, which is not present in the distractor grids. Furthermore, I will also consider the spatial arrangement of X's in the target grid, such as the fact that the X's in the first and last columns are aligned with the X's in the middle column, which is not the case in the distractor grids.} \\
            }
        }
    }
    & & \\ \\

    \theutterance \stepcounter{utterance}  
    & & & \multicolumn{4}{p{0.6\linewidth}}{
        \cellcolor[rgb]{0.95,0.95,0.95}{
            \makecell[{{p{\linewidth}}}]{
                \texttt{\tiny{[P1$\rightarrow$P2]}}
                \texttt{Post{-}episode Reflexion (not sent to GM):} \\
\texttt{New plan: I failed to match the given expression because I didn't carefully analyze the description. The expression states that the grid has X's in the first and last columns of each row, but I selected the third grid, which has X's in the first and last columns of each row, but not in the first and last columns of each row as the expression states. The correct interpretation is that the first and last X in each row are in the same column, not necessarily the first and last columns. However, I also failed to notice that the expression states "two X's in the middle three columns", which is not present in the third grid. Furthermore, I failed to notice that the expression states "first and last columns", which implies that the first and last X in each row are in the first and last columns, not in the middle column. Next time, I will carefully read the expression and look for a grid that matches the description. I will pay attention to the position of the X's in each row, the column in which the first and last X appear, and the presence of X's in the middle three columns. I will also pay attention to the fact that the expression states "first and last columns", which implies that the first and last X in each row are in the first and last columns, not in the middle column.} \\
            }
        }
    }
     \\ \\

\end{supertabular}
}

\end{document}
